\section{Causal Validation}
\label{sec:validation}

All validation machinery is standard; we specify our configuration and one
computational shortcut.

\paragraph{Activation patching.} For rule $R$ with pair
$(\mathbf{x}_{\text{base}}, \mathbf{x}_{\text{source}})$, we sweep residual-stream
activations over every layer $\ell$ and token position $p$, patch
$\mathbf{h}^{(\ell,p)}_{\text{base}} \leftarrow
\mathbf{h}^{(\ell,p)}_{\text{source}}$, and report the normalized causal
indirect effect on the logit difference $\Delta$ at the first divergent
answer token,
$\mathrm{CIE}_{\mathrm{norm}}(\ell,p) =
(\Delta_{\text{patched}} - \Delta_{\text{base}}) /
(\Delta_{\text{source}} - \Delta_{\text{base}})$,
where $1$ is full causal recovery; the normalization is not capped, so
overshoot can exceed $1$ \citep{vig2020,meng2022,zhang2024}.

\paragraph{Subspace interventions.} To localize each rule below the level
of a full residual-stream patch, we use interchange interventions on a
rank-$k$ orthonormal subspace ($k \in \{1,2,4,8\}$): the standard low-rank
DAS parameterization \citep{geiger2024das,wu2023boundless}, as implemented
in pyvene \citep{wu2024pyvene} and used by ReFT \citep{wu2024reft} and
AxBench \citep{wu2025axbench}. We claim no methodological novelty here.
Our one deviation is computational: instead of optimizing the interchange
objective through the frozen network, we fit the basis $W$ against a local
linear logit surrogate (ridge regression on counterfactual activation
pairs) and then validate the resulting subspace by causal intervention on
the original model (Appendix~\ref{app:subspace}). \TODO{report agreement
between surrogate-fitted and interchange-trained subspaces at matched
sites and ranks; if agreement is high, the surrogate is a cheap
approximation, otherwise we fall back to standard training}

\paragraph{Circuits and composition.} Path patching
\citep{wang2022ioi,goldowskydill2023} traces which attention heads
transmit each rule's subspace signal, retaining edges with normalized
effect $\ge 0.15$. For composition, each pair of independent rules
$(A,B)$ is tested with a four-prompt bundle (base, source $A$, source $B$,
joint $A{+}B$): we apply interventions in both orders and measure
order-sensitivity by KL divergence between output distributions, plus
recovery of the joint intervention (Appendix~\ref{app:commutativity}).

\section{Experimental Setup}
\label{sec:setup}

We evaluate four questions from three UKLO puzzles: 2023 Gilbertese Q3,
2019 Ndebele Q1, and Tshiluba Q7/Q8,\footnote{Puzzle sources are linked in
Appendix~\ref{app:rules}; we study only questions \texttt{Llama-3.1-8B}
answers exactly correctly in context, so interventions probe a mechanism
that succeeds (see Limitations).} spanning Austronesian VOS morphology and
Bantu SVO concord systems. All experiments use
\texttt{meta-llama/Llama-3.1-8B} \citep{grattafiori2024} via NNsight/NDIF
\citep{fiottokaufman2025}: 238{,}975 patched forward evaluations, subspace
analyses at 145 causal sites, 5{,}069 path-patching traces, and 31
paired-rule intervention runs (16 distinct rule pairs). Gilbertese
supports $N{=}18$ counterfactual training pairs (ranks up to 8); Ndebele
and Tshiluba support $N{=}2$ (ranks capped at 2).
